
\documentclass[11pt,a4paper]{article}
\usepackage[margin=2.5cm]{geometry}
\usepackage[utf8]{inputenc}
\usepackage[T1]{fontenc}
\usepackage[french]{babel}
\usepackage{enumitem}
\usepackage{hyperref}


\begin{document}
	\pagestyle{empty}
	
	\noindent
	Friedly WOLI \\  % \hfill Airbus \\
	Toulouse, 31100  \\
	0767339290 \\
	friedlysedjro@gmail.com 
	
	\begin{flushright}
		
		Université de Toulouse \\
		118 route de Narbonne 31062 
	\end{flushright}
	
	
	\bigskip
	
	
	
	
	
	
	\textbf{Objet :} Candidature master MSME \\ \\
	Madame, Monsieur;\\
	
	Actuellement étudiant en mécanique énergétique à l’Université de Toulouse, je souhaite intégrer le Master 1 Modélisation et Simulation en Mécanique et Énergétique (MSME) afin d’approfondir mes compétences en modélisation physique et en simulation numérique appliquées aux problématiques industrielles et de recherche.\\
	
	Au cours de mon parcours, j’ai acquis de solides bases en mécanique des fluides, en transferts thermiques et en mécanique des structures. J’ai développé un intérêt particulier pour la modélisation des phénomènes physiques complexes ainsi que pour leur résolution numérique. Les méthodes de calcul scientifique, la programmation notamment en Python ainsi que l’utilisation d’outils numériques m’ont permis de mieux appréhender le lien entre les approches théoriques et leurs applications industrielles.\\
	
	Le parcours MSME représente pour moi une continuité cohérente et exigeante. La place importante accordée aux méthodologies de simulation, aux projets numériques, ainsi qu’à l’utilisation d’outils professionnels tels qu’OpenFOAM, Abaqus ou COMSOL correspond pleinement à mon projet professionnel. Je suis particulièrement motivé par l’approche orientée vers la résolution de problématiques concrètes en mécanique énergétique, ainsi que par l’accent mis sur l’autonomie, la rigueur scientifique et la gestion de projet.\\
	
	Mon objectif professionnel est d’évoluer vers un poste d’ingénieur en bureau d’études ou en recherche et développement dans le secteur aéronautique. À plus long terme, je souhaite également me laisser la possibilité de poursuivre en doctorat afin de contribuer à des projets de recherche en modélisation avancée.\\
	
	Rigoureux, organisé et déterminé, je suis prêt à m’investir pleinement dans cette formation exigeante. J’apprécie particulièrement le travail par projet, qui permet de développer à la fois les compétences techniques, l’autonomie et la capacité à travailler en équipe.\\
	
	Je vous remercie par avance de l’attention que vous porterez à ma candidature et me tiens à votre disposition pour tout entretien.\\
	
	Je vous prie d’agréer, Madame, Monsieur, l’expression de mes salutations distinguées.
	
	
	
	\vspace{0.8cm}
	
	
	
	
	
	
	\begin{flushright}
		
		\textbf{Friedly WOLI}
	\end{flushright}
	
	
\end{document}




